\runninghead{The Three Interpretations Compared}
\section{The Three Interpretations Compared}
\label{sec:comparison}

\begin{comparisontable}{}{Both grow until the harvest}{Thoughts of peace, cries from the depths}{While men slept}
Question &
Why does the householder forbid the weeding, and what does that ask of those
who hear it? &
Why does God, who has promised, command his people to call on him? &
If the enemy sows while men sleep, how is the field kept? \\
Starting text &
The Gospel's \latin{ne forte \dots\ Sínite}, Mt 13:29--30 &
The Introit, Jer 29:11--12, 14, with the Alleluia and Offertory, Ps 129:1--2 &
The Gospel's \latin{cum dormírent hómines}, Mt 13:25, with the Epistle,
Col 3:15--16 \\
Principal authors &
Augustine, Jerome, Chrysostom, Rabanus &
Jerome, Augustine, Hilary, Chrysostom &
Chrysostom, Aquinas, Hilary \\
Its answer &
Those who judge may be deceived and the cockle may become wheat, so the
separation is left to the angels at the harvest &
God could fulfil his promise without prayer and commands it so that his
people may receive it; faith asks to be heard &
Moral sleep can be avoided; the word dwelling richly and the heart that
prays and sings keep the field \\
What it asks &
Judge slowly, seek conversion, bear injuries, forgive &
Pray within the promise, as sinners in the deep &
Keep the heart awake with the word \\
\end{comparisontable}

The three interpretations share two things. Each places the present before a
fulfilment not yet given: the harvest in the first and third, where judgment
waits for it and the watch is kept until it, and the return in the second.
Each finds in the Epistle the conduct of the present time.

They differ in what they ask the Gospel. The first and third use the same
verses of the parable, but different words of it: the first the servants'
question and the householder's answer, the third the sleepers whom the
answer never mentions. They also look at different objects. The first is
about how the members of the Church treat those they take for cockle, and its
patience is directed at others. The third is about the field of each hearer's
own heart, and its vigilance is directed at oneself. The second alone takes
the chants as its starting point, and uses the Gospel only for the term the
householder sets, until the harvest.

The second and third both use Psalm~129 and both use St Hilary on it, but
different sentences of his. The second takes \latin{fidei nostrae est ut
precemur audiri}, faith's part is to pray to be heard; the third takes
\latin{non labiis, sed secreto \dots\ cordis}, not with the lips but in the
secret of the heart. The first uses the same psalm only for St Augustine's
answer that the one who cries from the deep is a sinner, which puts the
servants who would weed among those who need mercy.

None of the three needs another's conclusion, and none contradicts another,
so a reader may hold all three at once. The one real disagreement among the
Fathers cited here lies inside the first interpretation. It concerns how far
the Church may correct those she takes for cockle before the harvest. Jerome
reserves the doubtful case to God and does not exclude putting away the
manifest one; Chrysostom forbids killing heretics and allows restraint short
of it; Augustine came to accept exile and fines while holding that the
separation belongs to the end; and Aquinas holds that heretics deserve
death, because corrupting the faith is graver than forging money, and that
the Church, after a first and second admonition, excommunicates the stubborn
heretic and delivers him to the secular tribunal for death. He answers the
parable by reading its command, on Augustine's authority, as applying where
the cockle cannot be removed without the wheat.
The disagreement does not touch the other two interpretations, and it does
not touch the ground of the first, which all four hold: the safety of the
wheat and the possible conversion of the cockle.
